% Complete arbitrary-framework proofs for P18-A1.
\section{Arbitrary abstract argumentation: proofs}
\label{sec:annex-p18-arbitrary-argumentation}

\begin{annexitem}{annex:p18.arbitrary.admissibility}{thm:p18-arbitrary-admissibility}{The full arbitrary-set existence and relation arguments}
\begin{proof}
If \(S\subseteq T\), every target of an attack from \(S\)
is also a target of an attack from \(T\). Thus
\(S^+\subseteq T^+\), and the defining universal
defense condition gives \(\mathsf F(S)\subseteq\mathsf F(T)\).
Suppose \(S\) is conflict-free and \(a,b\in\mathsf F(S)\)
with \(a\to b\). Defense of \(b\) gives \(s\in S\)
attacking \(a\); defense of \(a\) gives \(t\in S\)
attacking \(s\), a conflict in \(S\). This includes
the case \(a=b\). Hence \(\mathsf F(S)\) is
conflict-free. If \(S\subseteq\mathsf F(S)\),
monotonicity gives
\(\mathsf F(S)\subseteq\mathsf F^2(S)\), proving
admissibility of the image. Also
\[
 S\cup\{a\}\subseteq\mathsf F(S)
                   \subseteq\mathsf F(S\cup\{a\})
                \qquad(a\in\mathsf F(S)).
\]
The containing set \(\mathsf F(S)\) is conflict-free,
so this proves admissible enlargement. The same
monotonicity retains every other defended argument
\(a'\). No finite enumeration entered these proofs.

Let \(\mathcal D\) be nonempty and directed, with
each member admissible, and set \(U=\bigcup\mathcal D\).
If two elements of \(U\) formed an attack, take
members of \(\mathcal D\) containing each; a
common containing member would contain that attack,
contradicting its conflict-freeness. If \(a\in U\),
choose a member \(D\) containing \(a\).
Every attacker of \(a\) is already attacked by
\(D\subseteq U\). Thus \(U\) is admissible.
It is the least upper bound because every upper
bound contains its union. The empty set is
admissible and contained in every admissible set.

Fix an admissible seed \(S\). The poset of
admissible supersets of \(S\) is nonempty.
Every nonempty chain has its admissible union
as an upper bound; an empty chain has \(S\)
as an upper bound. Zorn's lemma therefore
gives a maximal member \(E\). Any admissible
proper superset of \(E\) would still contain
\(S\), so \(E\) is preferred in the full
framework. Taking \(S=\varnothing\) gives
preferred existence.

If a preferred \(E\) defended \(a\notin E\),
admissible enlargement would contradict maximality.
Thus \(\mathsf F(E)\subseteq E\); admissibility
gives the reverse inclusion and conflict-freeness.
In general the definition of a complete extension
is exactly these two inclusions and conflict-freeness,
proving its fixed-point characterization.

Let \(E\) be stable. An attacker \(b\) of
a member \(a\in E\) is outside \(E\), by
conflict-freeness, and is therefore attacked
by \(E\). This proves admissibility.
Every proper superset of \(E\) contains
an outsider attacked by \(E\), so is
not conflict-free. Hence \(E\) is
preferred. Finally conflict-freeness
gives \(E\subseteq A\setminus E^+\);
attacking every outsider gives the
opposite inclusion. Conversely the
equation \(E=A\setminus E^+\)
gives both properties directly.
For \(A=\varnothing\), its sole
subset is empty and all these
extension definitions hold for it.
\end{proof}
\end{annexitem}

\begin{annexitem}{annex:p18.transfinite.grounded}{thm:p18-transfinite-grounded}{Increasing stages, stabilization and leastness}
\begin{proof}
At stage zero the seed is admissible.
At a successor stage the previous theorem
preserves admissibility, and the defining
inclusion of an admissible set gives
\(S_\alpha\subseteq S_{\alpha+1}\).
At a nonzero limit stage, the earlier
stages form a nonempty chain. Its
union is admissible by the directed-union
result and contains every earlier
stage. These statements prove by
transfinite induction that all stages
are admissible and increasing.
Transfinite recursion is performed
only on a set-sized ordinal interval;
the following bound ensures that
such an interval suffices.

For infinite \(|A|=\kappa\), fix
a well-order of \(A\). If every
inclusion \(S_\alpha\subseteq S_{\alpha+1}\)
were strict for \(\alpha<\kappa^+\),
choose the first element of
\(S_{\alpha+1}\setminus S_\alpha\)
in that well-order. Distinct stages
give distinct elements, since the
sequence is increasing. This would
inject \(\kappa^+\) into \(A\),
contradicting its cardinality.
An equality therefore occurs
before \(\kappa^+\).
For finite \(A\), \(n+1\)
strict inclusions would require
more than \(n\) distinct elements;
thus an equality occurs by
stage \(n\). Once equality
occurs the defining successor
and limit clauses retain that
same set at every later stage.

Let \(T\) be any fixed point
with \(S\subseteq T\); do not
assume \(T\) conflict-free.
By transfinite induction,
\(S_\alpha\subseteq T\).
The successor step is
\(\mathsf F(S_\alpha)\subseteq
  \mathsf F(T)=T\), and
the limit step is a union
of subsets of \(T\).
The stabilized stage is
therefore the least fixed
point containing \(S\).
It is admissible and fixed,
so is complete.

For the empty seed every
fixed point contains the
seed, giving absolute
leastness and uniqueness.
This grounded fixed point
is itself a complete
extension and is contained
in all of them, so their
intersection equals it.
Extensiveness follows from
the increasing sequence.
If \(S\subseteq T\) are
admissible, then
\(\operatorname{Cl}(T)\)
is a fixed point containing
\(S\); leastness gives
\(\operatorname{Cl}(S)\subseteq
  \operatorname{Cl}(T)\).
Starting at a complete
extension is already
starting at a fixed point,
which proves idempotence.
\end{proof}
\end{annexitem}

\begin{annexitem}{annex:p18.complete.semilattice}{thm:p18-complete-semilattice}{All meets and the exact upper-bound construction}
\begin{proof}
Let \(\mathcal E\ne\varnothing\) consist
of complete extensions, let \(\mathcal L\)
be their admissible lower bounds, and
put \(M=\bigcup\mathcal L\).
At least the empty set lies in
\(\mathcal L\). Choose any \(E_0\in\mathcal E\).
The entire union \(M\) is contained
in the conflict-free set \(E_0\),
so is conflict-free, even though
\(\mathcal L\) need not be directed.
Each element of \(M\) is defended
inside one of its admissible
constituents, hence is defended
by \(M\). Thus \(M\) is admissible
and contained in every member
of \(\mathcal E\).

Its admissible image \(\mathsf F(M)\)
is also contained in each such
member, since
\(\mathsf F(M)\subseteq\mathsf F(E)=E\).
Therefore \(\mathsf F(M)\in\mathcal L\),
so \(\mathsf F(M)\subseteq M\).
Admissibility gives \(M\subseteq\mathsf F(M)\).
Consequently \(M\) is complete.
Every complete lower bound is
an admissible lower bound and
is contained in \(M\). This
proves the displayed meet formula.

For a nonempty directed family
\(\mathcal D\), its union \(U\)
is admissible by the previous
theorem. Its complete closure
contains every member. Any
complete upper bound contains
\(U\), and hence contains
\(\operatorname{Cl}(U)\)
by leastness. This proves
the directed supremum formula.
The grounded extension is
least among all complete
extensions, giving the
empty supremum.

For an arbitrary nonempty
family \(\mathcal E\), if
its union \(U\) is conflict-free,
it is admissible: an element's
defense is supplied inside
the complete extension from
which it came. Its closure
is consequently the same
least complete upper bound.
Conversely a common complete
upper bound contains \(U\)
and is conflict-free, so
\(U\) must be conflict-free.

The finite mutual-pair example
has complete singletons
\(\{a\},\{b\}\) whose union
conflicts. They have no
common upper bound, so the
poset need not be a
complete lattice. In the
shared-argument example, the
intersection of complete
\(\{a,d\}\) and \(\{b,d\}\)
is \(\{d\}\). The previously
computed admissible list
shows that its only
admissible subset is empty.
Their meet is therefore
empty, not their intersection.

For a directed union that
is not complete, use arguments
\(u_n,v_n\), \(n<\omega\),
and \(z\), with attacks
\(u_n\to v_n,v_n\to u_n,v_n\to z\).
For every finite \(m\),
\(E_m=\{u_n:n<m\}\) is
conflict-free and fixed:
it defends exactly its
chosen \(u_n\), no \(v_n\),
and not \(z\), whose
infinitely many attackers
have not all been attacked.
The union is \(U=\{u_n:n<\omega\}\).
It defends \(z\), so
\(\mathsf F(U)=U\cup\{z\}\ne U\).
The latter set is fixed
and conflict-free and hence
is the least complete
upper bound. This proves
the stated distinction
for directed suprema.
\end{proof}
\end{annexitem}

\begin{annexitem}{annex:p18.finitary.continuity}{thm:p18-finitary-continuity}{Finite attacker sets and directed-union continuity}
\begin{proof}
Let \(\mathcal D\) be a nonempty
directed family of arbitrary subsets
of \(A\), and let \(U=\bigcup\mathcal D\).
Monotonicity gives
\(\bigcup_{D\in\mathcal D}\mathsf F(D)
 \subseteq\mathsf F(U)\).
If \(a\in\mathsf F(U)\), list its
finitely many attackers. For each
attacker choose a defender from
\(U\), then a member of
\(\mathcal D\) containing it.
Directedness, applied finitely
many times, supplies a single
member containing all these
defenders. It defends \(a\).
If \(a\) has no attackers,
choose any member of the
nonempty family; it defends
\(a\) vacuously. This proves
the reverse inclusion.

For an admissible seed,
the finite iterates are
increasing. Their union
\(S_\omega\) therefore satisfies
\[
 \mathsf F(S_\omega)
   =\bigcup_{n<\omega}\mathsf F(S_n)
   =\bigcup_{n<\omega}S_{n+1}
   =S_\omega.
\]
Every fixed point containing
the seed contains every finite
iterate and thus their union.
It follows that this union
is exactly the complete closure.
The finite-stage and converse
limitations follow from the
explicit examples proved below.
Nonempty directed families
are essential to the statement:
\(\mathsf F(\varnothing)\)
may contain unattacked arguments,
so preservation of the
empty union was not asserted.
For a directed family of complete
extensions, its union \(U\)
is admissible and hence
conflict-free. Continuity
now gives
\(\mathsf F(U)=\bigcup_{E\in\mathcal D}\mathsf F(E)
             =\bigcup\mathcal D=U\).
Thus it is complete
and is the plain-union
least upper bound.
\end{proof}
\end{annexitem}

\begin{annexitem}{annex:p18.alternating.levels}{thm:p18-alternating-levels}{The complete alternating-level calculation}
\begin{proof}
If \(S\subseteq T\), then \(S^+\subseteq T^+\);
taking complements reverses that inclusion.
For any \(a\in A\),
\[
\begin{aligned}
 a\in\mathsf N(\mathsf N(S))
 &\ \Longleftrightarrow\
       \text{no attacker of }a\text{ belongs to }\mathsf N(S)\\
 &\ \Longleftrightarrow\
       \text{every attacker of }a\text{ belongs to }S^+\\
 &\ \Longleftrightarrow\ a\in\mathsf F(S).
\end{aligned}
\]
This includes an argument with no attackers,
since both universal clauses are then vacuous.
Thus two successive level operations equal
one characteristic operation.

The initial even level is \(A\).
The first odd level \(\mathsf N(A)\)
consists exactly of the unattacked
arguments, and therefore equals
\(\mathsf F(\varnothing)\).
Applying \(\mathsf N^2=\mathsf F\)
proves the two displayed level
formulas by induction.
The iterates from the empty
set increase and are contained
in the least fixed point \(G\).
The iterates from \(A\) decrease:
\(\mathsf F(A)\subseteq A\),
and monotonicity preserves this
inclusion at every subsequent
iterate. Since \(G\subseteq A\)
and \(\mathsf F(G)=G\),
monotonicity also gives
\(G\subseteq\mathsf F^n(A)\)
for every \(n\). This proves
all odd/even bounds.

If \(a\) is eventually in
every level, choose an odd
level beyond the specified
threshold. It lies in the
union of odd levels.
Conversely, if \(a\) belongs
to one odd level, it
belongs to every later
odd level by their
increase, and to every
even level because that
odd level is contained
in \(G\), itself contained
in each even level.
Thus it belongs to
all sufficiently late
levels. The indefeasible
set is exactly the
odd-level union.

Their formula identifies
this union with all
finite grounded iterates;
including the zero
iterate adds nothing.
It is contained in
\(G\). Under finitarity
the continuity theorem
identifies that union
with \(G\) itself.
The already specified
\(\omega+1\) example,
computed in the examples
proof below, has \(z\)
in \(G\) but in no
finite grounded iterate.
It therefore proves
the possible strictness.
\end{proof}
\end{annexitem}

\begin{annexitem}{annex:p18.residual.lifting}{thm:p18-residual-lifting}{All cross-boundary attackers and complete lifting}
\begin{proof}
Write \(B=B_E\). By its definition,
no member of \(E\) attacks \(B\).
If \(b\in B\) attacked \(e\in E\),
admissibility of \(E\) would supply
an attack from \(E\) to \(b\),
contradicting \(b\notin E^+\).
Thus no attack goes in the
other direction either.

For \(a\in B\), completeness
gives \(a\notin\mathsf F(E)\).
Some attacker \(b\) of \(a\)
is therefore not attacked
by \(E\). It cannot itself
belong to \(E\), since
then \(a\in E^+\).
Hence \(b\in B\), proving
the residual attacker assertion.

Let \(S\subseteq B\) be
admissible in the residual.
Both pieces \(E,S\) are
conflict-free and the
preceding paragraph excludes
all attacks between them.
For members of \(E\), defense
is already supplied by \(E\).
Consider an attacker \(b\)
of a member of \(S\).
If \(b\in E\), the attack
is impossible. If \(b\in E^+\),
it is already attacked by
\(E\). Otherwise \(b\in B\)
and residual admissibility
supplies a defender in \(S\).
This covers all attackers
and proves admissibility
of \(E\cup S\).

Now suppose \(S\) is complete
in the residual and \(a\)
is defended by \(E\cup S\).
If \(a\in E\), it is
already in the union.
If \(a\in E^+\), some
\(e\in E\) attacks \(a\).
Defending \(a\) would
require attacking that
\(e\) from the conflict-free
set \(E\cup S\) containing
it, a contradiction.
The remaining case is
\(a\in B\). Every residual
attacker of \(a\) is
attacked by \(E\cup S\);
it cannot be attacked
by \(E\), by the
definition of \(B\).
Thus \(S\) defends \(a\)
in the residual, and
its completeness gives
\(a\in S\). This proves
completeness of the union.
Finally \(B=\varnothing\)
means that every outsider
of the conflict-free \(E\)
is attacked by \(E\),
which is exactly stability.
\end{proof}
\end{annexitem}

\begin{annexitem}{annex:p18.well.founded.extensions}{thm:p18-well-founded-extensions}{Residual descent and uniqueness}
\begin{proof}
Take the grounded complete
extension \(G\). If it were
not stable, its residual
\(B_G\) would be nonempty.
Every member of that
residual has an attacker
in the residual. Starting
with any member, dependent
choice gives an infinite
sequence \(a_0,a_1,\ldots\)
there with \(a_{n+1}\to a_n\).
This contradicts well-foundedness.
Thus \(G\) is stable.

Every complete extension
contains \(G\). But a
stable set has no proper
conflict-free superset:
it attacks every new
outsider. Therefore every
complete extension equals
\(G\). It is preferred
by the stable-implies-preferred
result, and it is grounded
by construction. For the
empty framework the same
proof has empty residual
from the outset and
gives the sole extension
\(\varnothing\).
\end{proof}
\end{annexitem}

\begin{annexitem}{annex:p18.noncontroversial.seed}{thm:p18-noncontroversial-seed}{The full threat-defense iteration without a minimality premise}
\begin{proof}
Fix the seed \(a\). We first prove
that every member of every \(E_n\)
has an even attack walk to \(a\).
At stage zero use the length-zero
walk at \(a\). If \(d\in D_n\),
some threat \(b\) satisfies
\(d\to b\to e\) for \(e\in E_n\).
Appending that two-edge walk to
the even walk from \(e\) to
\(a\) gives the required even
walk from \(d\). Old members
retain their existing walks.
This is an ordinary finite
induction and permits repeated
vertices exactly as in the
source's definition.

If \(x,y\in E_n\) and \(x\to y\),
the even walk from \(y\) to
\(a\), preceded by that attack,
is odd. The even walk from
\(x\) to \(a\) still exists.
Thus \(C(x,a)\), contrary
to the seed premise.
This also excludes \(x=y\),
and proves conflict-freeness
of every stage. Their union
is conflict-free: two proposed
conflicting members already
lie together at the larger
of their two finite stages.

Take \(e\in E_\omega\), choose
\(n\) with \(e\in E_n\), and
let \(b\to e\). If \(E_n\)
already attacks \(b\), the
union does too. Otherwise
\(b\in T_n\). By the
hypothesis that every argument
has an attacker, some
\(d\to b\) exists, and all
such \(d\) belong to \(D_n\).
Hence \(E_{n+1}\) attacks
\(b\). This proves defense
of every member of the
union against every attacker.
The union is therefore
admissible, is nonempty,
and contains \(a\).
Its transfinite complete
closure supplies the stated
complete extension.

Now assume a nonempty limited
controversial framework.
If its grounded extension
is nonempty, use that
complete extension.
Otherwise \(\mathsf F(\varnothing)
\subseteq G=\varnothing\),
so every argument has
an attacker. There exists
some \(a\) with no incoming
\(C\) relation: if every
argument had such an
incoming relation, dependent
choice on the nonempty
set would produce the
prohibited infinite sequence.
Apply the construction just
proved at this \(a\).
For an empty framework no
nonempty subset exists,
so the explicitly stated
nonemptiness hypothesis
cannot be removed.
\end{proof}
\end{annexitem}

\begin{annexitem}{annex:p18.coherence}{thm:p18-coherence}{Inheritance, residual enlargement and stable existence}
\begin{proof}
An induced subframework of a
limited controversial framework
is limited controversial.
Indeed, any finite walk in
the induced framework is
also a walk in the original.
Thus its controversy relation
is a subrelation of the
original \(C\); an infinite
sequence in the former would
be one in the latter.
The same reasoning proves
inheritance of uncontroversiality.

Choose any preferred extension
\(E\), whose existence was
proved using admissible chains.
It is complete. If it
were not stable, its residual
would be nonempty and
limited controversial.
The preceding lemma supplies
a nonempty complete extension
\(S\) of that residual.
Residual lifting makes
\(E\cup S\) admissible,
indeed complete, in the
original framework.
Since \(S\subseteq B_E\)
is nonempty and disjoint
from \(E\), this is a
proper admissible superset
of \(E\), a contradiction.
Every preferred extension
is therefore stable.
Existence of a preferred
extension now gives existence
of a stable one.

If the framework is
uncontroversial, \(C\) is
empty and cannot support
an infinite backward
sequence. It is thus
limited controversial.
The empty framework is
included: its sole preferred
set is empty and stable,
while the nonempty-residual
branch is never invoked.
\end{proof}
\end{annexitem}

\begin{annexitem}{annex:p18.relative.grounding}{thm:p18-relative-grounding}{The two residual extensions and intersection of preferred extensions}
\begin{proof}
Let \(G\) be grounded and
\(B=B_G\). The specified
argument \(a\) belongs
to \(B\). The residual
is uncontroversial by
inheritance, and every
one of its arguments
has an attacker in
that residual, by the
lifting lemma. In
particular no residual
argument is controversial
with respect to \(a\).
Run the full threat-defense
iteration in the residual,
starting at \(a\).
Its admissible union
has a complete residual
closure \(S_1\) containing
\(a\). The complete
extension \(E_1=G\cup S_1\)
of the original framework
therefore contains \(a\).

Choose a residual attacker
\(b\to a\), whose existence
was proved above. Apply
the same construction
starting at \(b\), giving
a complete residual
extension \(S_2\) containing
\(b\). Then \(E_2=G\cup S_2\)
is complete in the
original framework and
attacks \(a\). This
proves both required
alternatives with their
original attack orientation.

For relative grounding,
\(G\) is contained in
every preferred extension,
since each is complete.
Take any \(a\notin G\).
If \(a\in G^+\), every
preferred extension contains
an attacker of \(a\)
from \(G\), so by
conflict-freeness none
contains \(a\).
Otherwise \(a\in B_G\);
the constructed \(E_2\)
is admissible and attacks
\(a\). Extend it to
a preferred extension
using the arbitrary-set
maximality result. That
preferred extension cannot
contain \(a\), because
it already contains an
attacker of \(a\).
Thus no argument outside
\(G\) belongs to every
preferred extension.
This proves equality
with their intersection.
The argument also covers
the empty framework,
where both sides are
empty and there is
one preferred extension.
\end{proof}
\end{annexitem}

\begin{annexitem}{annex:p18.arbitrary.examples}{thm:p18-arbitrary-examples}{Complete calculations and the infinite separating cases}
\begin{proof}
\paragraph{Closure at \(\omega\).}
In \(\mathcal F_\omega\), \(g_0\) is
unattacked, \(g_{n+1}\)'s sole attacker
is \(h_n\), and \(h_n\)'s sole
attacker is \(g_n\). A set consisting
only of \(g\)'s attacks exactly the
correspondingly indexed \(h\)'s.
Consequently, for
\(G_m=\{g_n:n<m\}\),
\[
                    \mathsf F(G_m)=\{g_n:n<m+1\}.
\]
No \(h_n\) is defended by these
sets, since they attack no \(g_n\).
This proves the finite stages
by induction, including \(m=0\).
Their union contains every \(g_n\),
and its characteristic image
is itself. No finite stage
is fixed. Every vertex has
at most one attacker, proving
finitarity. A backward attack
sequence starting at a fixed
\(g_n\) or \(h_n\) decreases
the index after at most two
steps and terminates at
the unattacked \(g_0\).
Thus the framework is
well-founded.

\paragraph{Closure at \(\omega+1\).}
In the augmented framework,
the \(g,h\) calculations
are unchanged. The new
argument \(z\) is defended
by a \(g\)-only set
exactly when that set
contains every \(g_n\),
since \(z\)'s attackers
are all the \(h_n\),
and only \(g_n\) attacks
\(h_n\). Hence \(z\)
appears for the first
time in \(\mathsf F(G_\omega)\).
It attacks nothing, so
adding it does not
change the defense
of any other argument.
This yields the fixed
point \(G_\omega\cup\{z\}\).
It is conflict-free and
attacks all its outsiders,
the \(h_n\), so is stable.
The attacker set of \(z\)
is infinite. Moreover
\[
 \bigcup_{m<\omega}\mathsf F(G_m)=G_\omega
        \ne G_\omega\cup\{z\}=\mathsf F\!\left(\bigcup_{m<\omega}G_m\right),
\]
which is the promised
failure of \(\omega\)-continuity.
Any backward sequence from
\(z\) chooses a particular
\(h_n\) in its first step
and then terminates by
the preceding index argument.
Every other starting point
has the same property.
Thus this example is
still well-founded.

\paragraph{Acyclic but not well-founded.}
In \(\mathcal F_{\rm reverse}\),
each \(n\) is attacked
only by \(n+1\), and
only \(n+2\) attacks
that attacker. Thus for
every subset \(S\),
\[
                       \mathsf F(S)=\{n:n+2\in S\}.
\]
At a fixed point, membership
is constant along each
parity class. The only
fixed points are empty,
all even numbers, all
odd numbers and all
natural numbers. The
last is not conflict-free;
the first three are.
These are therefore
exactly the complete
extensions, with the
empty one grounded.
The two parity sets
are inclusion-maximal
among them and attack
every outsider through
its immediate successor.
They are stable and
preferred. There are
no further preferred
extensions, since every
preferred set is complete.
Every edge decreases
the natural number,
so there is no finite
directed cycle, but
\(0,1,2,\ldots\) is
an infinite backward
attack sequence.

\paragraph{Continuity without finitarity.}
In \(\mathcal F_{\rm shareddefender}\),
\(d\) is always defended
because it is unattacked.
No \(b_n\) can be defended:
its attacker \(d\) is
attacked by no argument.
The argument \(z\) is
defended exactly when
\(d\in S\), since \(d\)
is the sole attacker
of each \(b_n\).
This proves the stated
formula for \(\mathsf F(S)\).
For a nonempty union,
\(d\) belongs to the
union exactly when it
belongs to one of its
members, so the formula
preserves that union
of characteristic values.
In particular it preserves
all nonempty directed
unions, although \(z\)
has infinitely many
attackers.

\paragraph{Limited but nonzero controversy.}
The walks \(c\to a\)
and \(c\to b\to a\)
give both parities
from \(c\) to \(a\).
These are the only
pair of endpoints
with both parities:
the graph has only
the three displayed
edges and no directed
cycle. Thus its
controversy relation
is exactly \(\{(c,a)\}\).
It cannot support
an infinite backward
sequence. The only
unattacked argument is
\(c\); neither \(a\)
nor \(b\) can be
defended against \(c\),
which is never attacked.
Thus \(G=\{c\}\) is
fixed and stable.
Every complete extension
contains it, and no
proper superset is
conflict-free, proving
the asserted uniqueness.

\paragraph{Dung's initial framework.}
The sole unattacked
argument is \(i_2\).
The set \(\{i_2\}\)
attacks \(a\), thereby
defending \(i_1\).
It cannot defend \(a\),
whose attacker \(i_2\)
is never attacked.
Consequently its image
is \(\{i_1,i_2\}\),
and the latter set
is fixed. It is
conflict-free and
attacks its sole
outsider \(a\), so
is stable. Grounded
containment in every
complete extension
and stability again
prove uniqueness of
the complete extension,
and hence of preferred
and stable extensions.
The attacks \(i_1\to a\to i_1\)
form a two-cycle, which
can be repeated to give
an infinite backward attack
sequence. Hence this
coincidence does not
require well-foundedness.

\paragraph{No inclusion-minimal defense.}
Partition the arguments
of \(\mathcal F_{\rm nomindefense}\)
into two classes: put
\(a\) and \(x_{k,j}\)
with even \(j\) in
one class, and all
\(b_n\) and \(x_{k,j}\)
with odd \(j\) in
the other. Every
attack changes class.
The parity of every
walk is consequently
determined by its
endpoint classes.
No pair can admit
both parities, proving
uncontroversiality.
The argument \(a\)
has attacker \(b_0\);
each \(b_n\) has
attacker \(x_{n,0}\);
and each \(x_{k,j}\)
has attacker \(x_{k,j+1}\).
Thus every argument
has an attacker.

The threats to \(\{a\}\)
are exactly all \(b_n\).
Only arguments \(x_{k,0}\)
attack those threats,
and \(x_{k,0}\) attacks
exactly those with \(n\leq k\).
A defense \(D\) must
therefore contain \(x_{k,0}\)
for an unbounded set
of indices \(k\);
this condition is
also sufficient.
Removing any one
such index leaves
an unbounded set.
Any other member
of \(D\) contributes
nothing to attacking
the threats and
can also be removed.
Every defense thus
has a proper subset
which is still a
defense. None is
inclusion-minimal.
All premises of the
full threat-defense
construction nevertheless
hold, and its proved
admissible union and
complete closure apply.
More explicitly, its stages
starting at \(a\) satisfy
\[
 E_n=\{a\}\cup\{x_{k,2j}:k<\omega,\ j<n\}.
\]
At the initial stage all
\(b_r\), \(r<\omega\), are threats and their
attackers give every
\(x_{k,0}\). At each
later stage the as-yet
unattacked threats are
the arguments \(x_{k,2n-1}\);
their attackers add
exactly \(x_{k,2n}\).
This proves the formula
inductively. The union
has \(a\) and all
even-indexed \(x\)'s.
It attacks all \(b_n\)
through the \(x_{k,0}\)
and all odd-indexed
\(x\)'s through their
next even-indexed attacker.
It is conflict-free,
so is stable.

The two final order-structure
claims are exactly the
mutual-pair and shared-argument
calculations proved in
the complete-semilattice
annex item. They show
why the order-theoretic
operations must retain
their declared poset.
For the shared-argument
framework, the controversy
relation is exactly
\[
                   \{(a,c),(b,c),(a,d),(b,d)\}.
\]
The mutual attacks between
\(a,b\) contribute only
even extra lengths when
returning to the same
vertex; reaching \(c\)
directly or via the
other vertex contributes
both parities, and the
edge \(c\to d\) retains
both parities at \(d\).
There are no other
outgoing paths from
\(c,d\) that could
create a further pair.
No member of \(\{c,d\}\)
is a source in this
controversy relation,
so no two controversy
edges can compose.
In particular it has
no infinite backward
sequence. The framework
is limited controversial.
Its previously computed
preferred and stable
families coincide,
but their intersection
\(\{d\}\) differs from
the grounded empty set.
Thus limited controversy
by itself does not
imply relative grounding.
Finally, the alternating-level
theorem identifies the
indefeasible set of
\(\mathcal F_{\omega+1}\)
with its finite grounded
union, which the earlier
calculation gave as all
\(g_n\) without \(z\).
Its grounded extension
contains the additional
\(z\), proving the
last claimed strict
inclusion.
\end{proof}
\end{annexitem}

\begin{annexitem}{annex:p18.finite.graph.procedures}{thm:p18-finite-graph-procedures}{Finite enumeration and parity reachability}
\begin{proof}
With \(n\) listed arguments,
there are exactly \(2^n\)
candidate subsets. Each
candidate's attack targets,
conflicts and defended
arguments are calculated
by finite scans of the
listed relation. The
admissibility, fixed-point
and outsider-attack tests
are precisely the defining
conditions. Comparing each
retained admissible set
against the finite list
of its admissible proper
supersets removes exactly
the nonmaximal members.
Thus the procedure
terminates and returns
exactly each requested
extension family.
Checking the finite
preferred list against
the stable list is
exactly coherence.
The preferred list is
nonempty by the
existence theorem, so
its finite intersection
is well-defined; comparing
that intersection with
the computed grounded
set is exactly relative
grounding.
The grounded procedure
already has its proved
\(n+1\) operator-evaluation
bound. For the empty
argument set the subset
enumeration has its
single candidate
\(\varnothing\), giving
the correct singleton
families.

Induct on walk length
in the doubled graph.
A length-\(m\) attack
walk from \(b\) to
\(a\) lifts uniquely,
from starting parity
zero, to a walk
ending at \((a,m\bmod2)\).
Conversely projecting
a lifted walk onto
its first coordinates
gives the original
attack walk with that
same parity. The
zero-length case
reaches \((b,0)\)
and accounts for
even self-reachability.
Therefore the simultaneous
reachability test computes
exactly \(C(b,a)\).
A finite search marking
visited doubled vertices
terminates: there are
only \(2n\) such
vertices. It preserves
all reachable vertices,
including those reached
by walks with repetitions.

In any finite directed
graph, a cycle gives
an infinite backward
walk by repeatedly
traversing its vertices
in predecessor order.
Conversely an infinite
backward walk repeats
a vertex because the
vertex set is finite.
The segment between
two repetitions is
a nonempty closed
directed walk; deleting
shorter repeated segments
gives a directed cycle.
Apply this equivalence
first to \(R\), then
to the computed \(C\).
It proves both
acyclicity criteria.
Uncontroversiality is
exactly emptiness of
\(C\) by definition.

No such repetition
principle holds for
an infinite set:
the explicit reverse
natural-number framework
has a backward sequence
with distinct vertices
and no cycle. Similarly,
the ordinal and countable
union constructions were
proved as operations
on sets and ordinal
stages; no finite
representation of those
sets or decision test
for their limit equality
was supplied. The
stated algorithmic
claims therefore have
exactly their finite
input domain.
\end{proof}
\end{annexitem}
